\section{Historial de Cambios y Justificación de Modificaciones}En este apartado se detallan las modificaciones realizadas entre la primera iteración del documento de requisitos y la segunda versión mejorada, justificando la motivación técnica y metodológica de cada cambio.\subsection{Modificaciones en la Definición de Personas y Escenarios}
\begin{itemize}
\item \textbf{Normalización de terminología e inclusión de roles inclusivos:}
\begin{itemize}
\item \textbf{Persona 1 (Familia):} Se añade el rol explícito de Madre en la ficha de Gema Uceda Cabezali y se modifica la sección Familia para alinearse con la visión inclusiva de la comunidad educativa.
\item \textbf{Persona 2 (Profesor/a):} Se actualiza la denominación de Profesor a Profesor/a en todo el documento y se reemplaza el término alumnos por estudiantes para mantener consistencia con el vocabulario inclusivo del proyecto.
\end{itemize}\item \textbf{Reescritura y enriquecimiento de escenarios reales:}
\begin{itemize}
    \item \textbf{Escenarios de la Familia (Persona 1):} En el Escenario 1 se especifica detalladamente la comprobación del calendario antes de añadir la tarea para evitar solapamientos y se amplía la descripción indicando la inclusión de un vídeo explicativo. Además, se añade un nuevo Escenario 2 focalizado en la eliminación de tareas inadecuadas para adaptar la agenda a la evolución del estudiante.
    \item \textbf{Escenarios de la Profesor/a (Persona 2):} Se incorpora un nuevo Escenario 2 centrado en la creación, gestión de datos personales y actualización del perfil educativo de un estudiante recién matriculado.
    \item \textbf{Escenario del Estudiante con Visibilidad Reducida (Persona 3):} Se profundiza en la interacción con el sistema en el Escenario 1 (añadiendo detalles como la visita al edificio de vida adulta, reprografía, aula de Vida Independiente y Sala de Luz Negra) y en el Escenario 2 (explicando la interacción mediante cuadrícula adaptada, nota de audio y botones gigantes para la hora).
    \item \textbf{Escenario de la Estudiante con Discapacidad Intelectual (Persona 5):} Se amplía el Escenario 1 describiendo el flujo de inicio de sesión accesible mediante contraseña de secuencia de pictogramas y la personalización por gama de colores en la interfaz.
\end{itemize}
\end{itemize}\subsection{Ajustes en el Plan de Entregas y Cálculo de Velocidad}
\begin{itemize}
\item \textbf{Replanificación de las Iteraciones y Velocidad:}
\begin{itemize}
\item Se reajusta el desglose de Puntos de Historia (PH) por iteración: la $1^{\text{a}}$ iteración pasa de 30 a 33 PH; la $3^{\text{a}}$ iteración de 32 a 31 PH; y se formalizan la $4^{\text{a}}$ iteración (8 PH) y la $5^{\text{a}}$ iteración (26 PH).
\item Se añade la equivalencia explícita del esfuerzo por punto de historia ($1 \text{ PH} \approx 4 \text{ horas}$, lo que representa 26.4 horas de trabajo individual por desarrollador por iteración).
\end{itemize}\item \textbf{Clarificación del alcance ampliado de las iteraciones:}
\begin{itemize}
    \item Se añaden formalmente al plan la Iteración 4 (orientada a funcionalidades de calidad de vida) e Iteración 5 (sistema de actividades optativas) para delimitar de manera explícita el trabajo dentro y fuera de alcance.
\end{itemize}
\end{itemize}\subsection{Actualización del Backlog del Producto y Tarjetas de Historias de Usuario (HU)}
\begin{itemize}
\item \textbf{Clarificación técnica de la creación de actividades (HU-P3 y HU-F7):}
\begin{itemize}
\item Se modifica el título y la descripción de las historias de usuario de creación de actividades tanto para el profesorado (\texttt{HU-P3}) como para las familias (\texttt{HU-F7}), haciendo explícito que se pueden crear actividades con pasos o sin pasos. Esto responde a la necesidad de permitir tareas simples y tareas guiadas secuencialmente.
\end{itemize}\item \textbf{Autenticación adaptada para estudiantes (\texttt{HU-E1}):}
\begin{itemize}
    \item Se actualiza la redacción para destacar que el inicio de sesión del estudiante se realiza en función de sus necesidades de accesibilidad, justificando el uso de pictogramas o métodos simplificados explicados en los escenarios.
\end{itemize}

\item \textbf{Rol de administración técnica (\texttt{HUT-A0}):}
\begin{itemize}
    \item Se redefine la tarjeta \texttt{HUT-A0} explicitando el rol de Administrador técnico de nuestra empresa, separando la gestión del entorno y despliegue del rol del usuario administrativo del centro educativo.
\end{itemize}

\item \textbf{Modificaciones en la asignación de iteraciones del Backlog:}
\begin{itemize}
    \item Se reasignan las historias \texttt{HU-E11} (aviso por cancelación) y \texttt{HU-F10} (canal de comunicación) excluyéndolas para centrar la entrega inicial estrictamente en las necesidades del Producto Mínimo Viable (PMV).
\end{itemize}
\end{itemize}